% ============================================
% SECTION 9: RECONSTRUCTION PRINCIPLES
% Building Cognitive Sovereignty Within Thermodynamic Constraints
% ============================================
%
% REVISION v2.1 NOTES:
% All v2 changes retained, plus:
% 1. Equation now REFERENCES Section 4 (no competing formulation)
% 2. Operationalized new constructs (Entropy Rate, Substitutability, etc.)
% 3. Added March (1991) citation for 20% slack / exploration-exploitation
% 4. Softened Chaotic domain AI language ("non-authoritative" not "noise")
% 5. Reframed existence proof as autoethnographic vignette + external cases
% 6. Fixed cross-reference bug (sec:equity not sec:measurement)
% 7. De-escalated final paragraph (removed pre-characterization of critics)
% 8. Added "Boundary Conditions" subsection
% 9. Added table footnote linking evaluation signals to Section 9.8
% 10. Qualified "AI capability thresholds" predictive claim
% 11. Ensured DeLaCroix2018 citation key consistency
% ============================================

\section{Reconstruction Principles: Building Cognitive Sovereignty Within Thermodynamic Constraints}
\label{sec:reconstruction}

\subsection{The Epistemological Prerequisites}

Before proposing reconstruction pathways, we must acknowledge a fundamental paradox: those seeking to rebuild cognitive sovereignty are themselves products of the vectorization they seek to escape. The map is never the territory \citep{korzybski1933}, and our maps were drawn with vectorized tools.

Contemporary decision science reveals the illusion of individual rationality. As \citet{kahneman2011} demonstrated, human decision-making operates through systematic biases rather than rational calculation. More fundamentally, distributed cognition theory \citep{hutchins1995} establishes that thinking occurs not within individual minds but across socio-technical systems. What we experience as ``our'' decisions emerge from complex interactions between cultural values, institutional practices, environmental constraints, and collective sense-making processes.

Highly formalized training can overfit people to legible problems and under-train them for domain ambiguity. The analytical tools revealing the vectorization problem are themselves products of vectorized thinking. We are attempting to examine the lens through which we see using that same lens: \citet{hofstadter1979} strange loop made manifest.

Critical retrospection of any recent non-trivial decision reveals this embedding. The choice to pursue alternative educational approaches, implement new organizational structures, or resist AI integration emerges not from individual cognition but from:
\begin{itemize}
\item Cultural layer: Prevailing beliefs about knowledge and value \citep{schein1985}
\item Institutional layer: Organizational structures and incentive systems \citep{dimaggio1983}
\item Social layer: Peer networks and professional communities \citep{granovetter1985}
\item Individual layer: Personal history and embodied experience \citep{bourdieu1990}
\end{itemize}

The individual functions as one node in a distributed processing network: ``the literal neuron of a bigger brain.'' Distributed does not mean powerless; it means leverage points exist at multiple scales. This recognition demands epistemic humility. We cannot stand outside the system to reconstruct it. As \citet{maturana1987} establish, we are ``structurally coupled'' to the environment we seek to transform.

\subsection{The Cynefin Diagnostic Framework}

Snowden's Cynefin framework \citep{snowden2007} provides the essential diagnostic tool for understanding where reconstruction is both necessary and possible. The framework distinguishes five domains, each requiring different cognitive architectures and intervention strategies.

\begin{table}[htbp]
\centering
\caption{Cynefin Reconstruction Design Matrix\textsuperscript{a}}
\label{tab:cynefin-matrix}
\small
\begin{tabular}{p{2cm}p{3cm}p{3cm}p{3.5cm}}
\hline
\textbf{Domain} & \textbf{Suitable Interventions} & \textbf{Failure Modes} & \textbf{Evaluation Signals} \\
\hline
Clear & Automate + standardize; best practices & Over-automation of edge cases; complacency & Error rate; throughput; drift monitoring \\
\hline
Complicated & Expert review + checklists; good practices & Expert bottlenecks; analysis paralysis & Peer audit accuracy; decision latency \\
\hline
Complex & Safe-to-fail probes + narrative capture; emergent practices & Premature convergence; probe-aversion & Learning velocity; resilience indicators; adaptation rate \\
\hline
Chaotic & Stabilizing action + after-action learning; novel practices & Paralysis; inappropriate deliberation & Time-to-stabilize; harm minimization \\
\hline
Confused & Escalation protocol + reclassification & Misclassification persistence; domain forcing & Time spent in disorder; reclassification accuracy \\
\hline
\end{tabular}

\vspace{0.5em}
{\footnotesize \textsuperscript{a}Evaluation signals are illustrative proxies; operational definitions appear in Section~\ref{sec:measurement}.}
\end{table}

The critical insight: organizations systematically misclassify complex problems as complicated, applying ``best practices'' where emergent approaches are needed. \citet{ford2024} quantified this misclassification: 341 ``Simple'' problems were actually Complicated, 437 ``Complicated'' were Complex, creating ``catastrophic failure'' when complicated approaches fail in complex domains.

Table~\ref{tab:cynefin-matrix} operationalizes the diagnostic framework into a decision instrument. The reconstruction challenge is not merely cognitive but organizational: building systems capable of accurate domain classification and appropriate response selection.

\subsection{Thermodynamic Constraints on Reconstruction}

Reconstruction faces inescapable thermodynamic constraints. The Second Law is a constraint on slack, maintenance, and reversibility---not a full causal explanation, but an analytic boundary that strategies must respect.

\subsubsection{The Sovereignty Budget Constraint}

The power-budget allocation model (Equation~\ref{eq:sovereignty}, Section~\ref{sec:methodology}) establishes the thermodynamic constraint: cognitive sovereignty requires sustained allocation of the brain's baseline power toward synthesis-supporting states, with architectural quality determining extraction resistance.

\begin{equation}
S = P_{brain} \times \alpha_{synth} \times R
\end{equation}

\noindent Where $P_{brain} \approx 20$W represents baseline cerebral metabolism, $\alpha_{synth} \in [0,1]$ represents synthesis-time allocation, and $R \in [0,1]$ represents extraction resistance. This is the framework's core constraint: Energy Investment Rate must exceed Entropy Rate, or cognitive architecture degrades toward the thermodynamic floor where patterns become simple enough for algorithmic extraction.

For reconstruction planning, these terms admit operational proxies:

\begin{itemize}
\item \textbf{Synthesis Allocation} ($\alpha_{synth}$): Hours of deep work + deliberate practice + embodied practice per week, as proxy for synthesis-state time fraction
\item \textbf{Extraction Resistance} ($R$): Operationalized via substitutability, transfer distance, tacit load, and integration density (defined below)
\item \textbf{Entropy Rate}: Rate of architectural degradation from extraction pressure
\end{itemize}

\subsubsection{Operationalizing the Proxies}

The following constructs require operational specification to avoid hand-waving:

\begin{itemize}
\item \textbf{Entropy Rate} (degradation pressure): Interruptions per day, context-switch frequency, sleep debt accumulation, administrative load fraction, and rework cycles. Higher values indicate faster architectural erosion \citep{mark2008, gonzalez2004}.

\item \textbf{Substitutability} (automation exposure): Degree to which tasks decompose into rule-following steps, availability of standardized input-output pairs, and existence of comparable automation benchmarks in the domain \citep{frey2017, autor2015}.

\item \textbf{Transfer distance}: Number of domain-boundary crossings required for an artifact or decision, assessed via domain-tag schemes. Higher transfer distance indicates greater integration demand.

\item \textbf{Tacit load} (explainability gap): Proportion of task performance that resists explicit articulation. Assessed via protocol analysis or apprenticeship duration requirements \citep{collins2010}.

\item \textbf{Integration density}: Network density of cross-references in personal knowledge systems, or count of cross-domain analogies deployed per unit time. Higher density indicates sphere-like architecture.
\end{itemize}

These remain proxies, not direct measures. They provide calibration for reconstruction efforts while acknowledging that cognitive architecture itself resists direct observation.

\subsubsection{Historical Investment Benchmarks}

Historical evidence reveals the investment scale required for different knowledge types:
\begin{itemize}
\item \textbf{Sophia} (theoretical wisdom): 20+ years sustained investment across multiple domains
\item \textbf{Phronesis} (practical wisdom): Lifetime daily practice with contextual feedback
\item \textbf{Techne} (craft knowledge): Decade-scale deliberate practice with embodied feedback loops; meta-analyses confirm substantial time investment though magnitude varies strongly by domain \citep{macnamara2014}, while historical apprenticeships required 7--10 years before journeyman recognition \citep{delacroix2018}
\item \textbf{Metis} (adaptive cunning): Constant exposure to novel challenges; cannot be scheduled
\item \textbf{Nous} (intuitive insight): Unknown threshold, possibly requiring conditions no longer reproducible
\end{itemize}

Modern constraints make full reconstruction rare; partial reconstruction requires deliberate slack engineering:
\begin{itemize}
\item Work demands: 40--60 hours/week of vectorized activity
\item Economic pressure: Continuous productivity requirements consuming baseline cognition
\item Attention economy: Optimized extraction mechanisms operating continuously \citep{zuboff2019}
\item Social expectations: Optimization for measurable outcomes at every scale
\end{itemize}

\subsection{The Equity Constraint}
\label{sec:equity}

Cognitive sovereignty is unevenly accessible. Reconstruction presumes access to slack: time, money, health, safety, community. Those lacking economic security cannot protect synthesis time; those lacking social support cannot sustain counter-cultural practice; those facing discrimination encounter additional cognitive load that consumes the very resources reconstruction requires.

This recognition has two implications. First, individual reconstruction protocols (Section~\ref{sec:individual-protocols}) must be understood as conditional on resource availability, not universally applicable self-improvement advice. Second, institutional reconstruction (Section~\ref{sec:organizational-architecture}) becomes not merely strategic but ethically necessary: organizations and societies must provision the slack that individuals cannot secure alone.

Reconstruction is not personal virtue; it is collective infrastructure. The thermodynamic requirements apply regardless of individual merit or motivation. Systems that fail to provide protected synthesis time, externalized survival load, and extended time horizons will produce vectors regardless of curriculum content, pedagogical innovation, or inspirational leadership.

\subsection{Boundary Conditions and Failure of Applicability}
\label{sec:boundary-conditions}

The reconstruction protocols that follow assume baseline stability. They fail under conditions including:

\begin{itemize}
\item \textbf{Acute crisis}: War, displacement, acute illness, or immediate survival threat. Cognitive resources necessarily redirect to threat response; synthesis becomes impossible and inappropriate.
\item \textbf{Extreme precarity}: Housing instability, food insecurity, or imminent financial collapse. The cognitive load of survival planning consumes all available slack.
\item \textbf{Severe caregiving load}: Primary responsibility for dependents with high needs. Available cognitive resources are legitimately allocated elsewhere.
\item \textbf{Institutional capture}: Employment conditions that permit zero discretionary time or penalize non-optimized activity. Exit may be prerequisite to reconstruction.
\item \textbf{Health constraints}: Chronic conditions affecting cognitive capacity, energy availability, or time horizons. Protocols require adaptation to actual capacity, not normative expectations.
\end{itemize}

These are not excuses but constraints. Acknowledging them transforms equity from ethical note into methodological boundary: the protocols apply where baseline conditions exist, and institutional intervention is required where they do not.

\subsection{Individual Reconstruction Protocols}
\label{sec:individual-protocols}

The following protocols assume access to baseline resources: stable housing, adequate nutrition, physical safety, and approximately 10 hours/week of discretionary time. Where these conditions are absent, individual reconstruction becomes impossible without prior institutional intervention (see Section~\ref{sec:equity}).

\subsubsection{Phase 1: Diagnostic Assessment (Month 1)}

\textbf{Current State Analysis}
\begin{itemize}
\item Map daily activities to Cynefin domains using Table~\ref{tab:cynefin-matrix}
\item Identify vector lock-in patterns (tasks that could be fully specified as algorithms)
\item Calculate actual discretionary time/energy after survival maintenance
\item Assess extraction vulnerability using substitutability indicators
\end{itemize}

\textbf{Sphere Potential Mapping}
\begin{itemize}
\item Identify existing cross-domain connections (latent integration capacity)
\item Recognize capacities from pre-vectorization periods (childhood interests, abandoned pursuits)
\item Map curiosity patterns that resist optimization pressure
\item Locate energy reserves available for investment
\end{itemize}

\textbf{Realistic Timeline Acceptance}

Accept thermodynamic reality:
\begin{itemize}
\item No shortcuts exist (constraints are non-negotiable)
\item Decade-scale investment required for substantive sphere development
\item Starting now provides meaningful time before labor-market exposure to AI substitution intensifies
\item Partial development offers meaningful protection; perfectionism is vectorized thinking
\end{itemize}

\subsubsection{Phase 2: Foundation Building (Months 2--6)}

\textbf{Cross-Domain Exposure}
\begin{itemize}
\item Read deliberately outside primary field (minimum 2 hours/week protected)
\item Attend events in unfamiliar domains without optimization goals
\item Engage communities with different knowledge traditions
\item Follow curiosity without demanding measurable outcomes
\end{itemize}

\textbf{Pattern Recognition Development}
\begin{itemize}
\item Maintain synthesis journal for cross-domain insights
\item Practice analogical thinking between disparate fields
\item Build personal pattern library (not filing system---integration network)
\item Resist premature categorization; tolerate ambiguity
\end{itemize}

\textbf{Embodied Practice Initiation}

Essential for developing non-extractable knowledge \citep{collins2010}:
\begin{itemize}
\item Physical craft (woodworking, pottery, gardening)
\item Movement practice (martial arts, dance, climbing)
\item Musical instrument learning
\item Somatic awareness development
\end{itemize}

\subsubsection{Phase 3: Cultivation (Months 6--24)}

\textbf{Deliberate Integration}
\begin{itemize}
\item Create projects requiring multiple domain knowledge
\item Write integrative analyses crossing disciplinary boundaries
\item Teach others using cross-domain metaphors (synthesis requires articulation)
\item Build integration as default cognitive mode
\end{itemize}

\textbf{Complexity Navigation Practice}
\begin{itemize}
\item Seek problems without predetermined solutions
\item Practice probe-sense-respond in safe contexts
\item Embrace emergence and uncertainty as signals, not failures
\item Document pattern recognition development for calibration
\end{itemize}

\textbf{Community Building}

Sphere development requires collective support---individual reconstruction is thermodynamically implausible:
\begin{itemize}
\item Find others attempting reconstruction (shared context reduces explanation cost)
\item Create learning partnerships with complementary domains
\item Share failures and insights (distributed error-correction)
\item Build counter-cultural spaces protecting synthesis time
\end{itemize}

\subsection{Organizational Reconstruction Architecture}
\label{sec:organizational-architecture}

\subsubsection{From Silos to Sphere Teams}

Replace functional specialization with cross-domain capacity:
\begin{itemize}
\item Pilot teams mixing 3--5 different expertise domains
\item Complex challenge focus rather than efficiency optimization
\item Protected learning time: minimum 20\% allocation to non-productive exploration. Below this threshold, exploration collapses into exploitation and learning becomes extraction \citep{march1991}---an organizational entropy tax that cannot be avoided.
\item Success measured by adaptation capacity, not standardization compliance
\end{itemize}

\subsubsection{From Best Practices to Emergent Practices}

Shift from standardization to contextual response:
\begin{itemize}
\item Train all staff in Cynefin framework and domain classification
\item Create ``probe budgets'' for complex challenges (resources allocated to safe-to-fail experiments)
\item Document emergent solutions as narratives, not procedures
\item Celebrate contextual wisdom over universal protocols
\end{itemize}

\subsubsection{From Extraction to Investment}

Reverse the energy flow:
\begin{itemize}
\item Implement genuine apprenticeship programs (3+ years, not bootcamps)
\item Create internal learning structures for cross-domain development
\item Protect thinking time from productivity metrics (unmeasured is not unproductive)
\item Measure knowledge development trajectories, not just application outputs
\end{itemize}

\subsection{Governance of Tools in AI-Saturated Environments}
\label{sec:governance}

Reconstruction does not require AI avoidance---an impossibility in contemporary environments. It requires governance structures that preserve human cognitive development while leveraging algorithmic capacity appropriately.

\subsubsection{Domain-Appropriate Deployment}

\begin{itemize}
\item \textbf{Clear domain}: AI deployment appropriate for throughput optimization; human oversight for edge-case detection
\item \textbf{Complicated domain}: AI as analysis support; human experts retain decision authority
\item \textbf{Complex domain}: AI prohibited as decision authority; may assist in pattern surfacing, option generation, and documentation, but humans retain probe-sense-respond agency
\item \textbf{Chaotic domain}: AI outputs are non-authoritative; use only as optional situational prompts while humans act to stabilize
\item \textbf{Confused domain}: AI classification assistance only; human meta-cognition required for domain determination
\end{itemize}

\subsubsection{Structural Safeguards}

\begin{itemize}
\item Require ``human appeal paths'' for any automated decision affecting rights, resources, or relationships
\item Treat all model outputs as hypotheses requiring independent verification
\item Track ``deskilling risk'' as first-class operational metric
\item Implement ``synthesis breaks''---periods of AI-free cognitive work to maintain integration capacity
\end{itemize}

\subsubsection{Practical Implementation}

Cynefin-informed AI governance can be operationalized within existing organizational constraints. Frameworks such as the author's ``Karten auf den Tisch'' protocol \citep{cotoaga2024karten} demonstrate that reconstruction tools can be developed and deployed without requiring complete system exit. The challenge is not technological but organizational: building governance structures that treat cognitive sovereignty as a first-class design constraint rather than an efficiency impediment.\footnote{The ``Karten auf den Tisch'' framework operationalizes domain classification and AI deployment boundaries for organizational decision-making. Available at \url{http://cotoaga.ai/karten-auf-den-tisch/}.}

\subsection{Measurement and Evaluation}
\label{sec:measurement}

Reconstruction requires outcome measures to distinguish from manifesto. The following indicators operationalize sovereignty development:

\subsubsection{Individual Indicators}

\begin{itemize}
\item \textbf{Complex work ratio}: Percentage of time spent in complex/chaotic domains without premature reduction to complicated/clear
\item \textbf{Integration artifact production}: Weekly cross-domain outputs (notes, diagrams, syntheses) demonstrating connection-building
\item \textbf{Embodied practice cadence}: Hours/week in somatic skill development
\item \textbf{Automation dependency index}: Frequency of inability to proceed without tool outputs; lower indicates higher sovereignty
\item \textbf{Transfer distance achieved}: Demonstrated application of learning across domain boundaries
\end{itemize}

\subsubsection{Organizational Indicators}

\begin{itemize}
\item \textbf{Probe ratio}: Fraction of initiatives run as safe-to-fail experiments (complex domain appropriate)
\item \textbf{Protected time utilization}: Percentage of allocated learning time actually taken (not ``available but unused'')
\item \textbf{Drift/incident rates}: Failures in automated decision systems requiring human intervention
\item \textbf{Apprenticeship metrics}: Retention rates and skill progression in long-term development programs
\item \textbf{Cross-domain mobility}: Internal movement between functional areas indicating integration capacity
\item \textbf{Reclassification accuracy}: Correct domain diagnosis in post-hoc analysis
\end{itemize}

These indicators are proxies, not direct sovereignty measures. They provide calibration for reconstruction efforts while acknowledging that cognitive architecture itself resists direct observation.

\subsection{Existence Proofs: Documented Reconstruction Trajectories}
\label{sec:existence-proof}

The thermodynamic analysis might suggest reconstruction is merely theoretical---possible in principle but unrealized in practice. Evidence to the contrary requires documented cases of vector trap escape under contemporary conditions.

\subsubsection{Methodological Note}

This section employs autoethnographic vignette as mechanism-illustration rather than representative evidence. The author's trajectory demonstrates \textit{feasibility} under specific (and atypical) resource conditions, not universal applicability. External cases provide triangulation.

\subsubsection{Vignette: Six Vector Trap Navigations}

The author's trajectory provides one existence proof. Six distinct vector traps encountered and navigated:

\begin{enumerate}
\item \textbf{Academic vector trap}: PhD program in quantitative finance (1990s). Exit: Abandoned for internet entrepreneurship when the field's vectorization trajectory became apparent.
\item \textbf{Technical vector trap}: Employee \#5 at Germany's first online brokerage; built technical chart analysis system (KAPITALIst). Exit: Recognized pattern-matching as ultimately extractable; pivoted to synthesis roles.
\item \textbf{Optimization vector trap}: Traveling Salesman Problem algorithm achieving global ranking leadership for six months (2004). Exit: Understood that algorithmic excellence without integration produces sophisticated vectors, not spheres.
\item \textbf{Geographic vector trap}: New York environment optimizing for extraction. Exit: Physical relocation (overland to South America) as literal escape from optimization pressure.
\item \textbf{Corporate vector trap}: Management consulting roles offering vectorized ``strategic'' work. Exit: Independent practice preserving synthesis capacity.
\item \textbf{Product vector trap}: Platform development pressure toward feature optimization. Exit: Maintained integration focus over standard product management vectorization.
\end{enumerate}

These escapes required: recognition of vector pressure before capture became irreversible, resources enabling transition (often accumulated during vector phases), willingness to sacrifice optimization rewards for sovereignty preservation, and extended time horizons incompatible with quarterly thinking.

\subsubsection{External Triangulation}

The autoethnographic vignette gains credibility through external cases demonstrating similar patterns:

\begin{itemize}
\item \textbf{Craft apprenticeship resurgence}: The revival of traditional apprenticeship in woodworking, blacksmithing, and textile arts documents intentional rejection of credentialized training in favor of embodied knowledge transmission \citep{sennett2008}.
\item \textbf{Slow science movement}: Academic communities explicitly resisting publication metrics and grant optimization in favor of long-term research programs demonstrate institutional-scale reconstruction attempts \citep{stengers2018}.
\item \textbf{Complexity practitioner networks}: Communities organized around Cynefin and related frameworks show sustained counter-practice against best-practice ideology in organizational consulting.
\end{itemize}

The existence proofs share common structure: recognition of vectorization pressure, resource accumulation enabling transition, deliberate sacrifice of optimization rewards, and community support for counter-cultural practice. The question shifts from ``Is reconstruction possible?'' to ``What conditions enable it?''---a question tractable through research rather than philosophy alone.

\subsection{The Budgetary Choice}

The choice is not binary; it is budgetary. Every system---individual, organizational, civilizational---allocates resources between sovereignty investment and extraction optimization, explicitly or implicitly.

\textbf{Full sovereignty investment} requires conditions approaching historical patronage: externalized survival costs, protected synthesis time, extended time horizons, and freedom from productivity measurement. These conditions existed for medieval scholars and Renaissance polymaths; they do not exist at civilizational scale today.

\textbf{Partial sovereignty investment} represents the realistic frontier: deliberate allocation of constrained resources toward reconstruction within extraction-dominated environments. This is not compromise but strategic necessity---the thermodynamically feasible path.

\textbf{Zero sovereignty investment}---whether through unawareness, resource absence, or optimization pressure---produces complete vectorization. This path offers immediate efficiency gains and guaranteed long-term obsolescence.

The reconstruction principles presented here assume partial investment: readers operating within extraction systems who seek to preserve or develop synthesis capacity. Full reconstruction would require institutional transformation beyond individual scope; zero reconstruction requires no guidance.

\subsection{Critical Success Factors and Failure Patterns}

\textbf{Success Requirements}:
\begin{itemize}
\item Accept decade-scale timeline (thermodynamic reality, not motivational failure)
\item Protect energy investment despite productivity pressure
\item Build community support (individual reconstruction is thermodynamically implausible)
\item Measure progress in capability development, not credential acquisition
\item Maintain sovereignty despite continuous extraction attempts
\end{itemize}

\textbf{Common Failure Patterns}:
\begin{itemize}
\item Treating reconstruction as temporary project rather than permanent practice
\item Attempting alone without community support (energy requirements exceed individual capacity)
\item Measuring by vectorized metrics (efficiency, speed, output volume)
\item Expecting linear progress in complex domain development
\item Underestimating energy investment required; abandoning after insufficient duration
\item Optimizing the reconstruction process itself (meta-vectorization)
\end{itemize}

\textbf{The Fundamental Recognition}:

Reconstruction is not optimization, self-improvement, or productivity enhancement. It is choosing cognitive sovereignty over efficiency, wisdom over productivity, and decades of investment over immediate returns. The choice is budgetary: allocate resources toward extraction resistance or accept dissolution into algorithmic substitutability.

A common objection holds that framing constraints as non-negotiable constitutes ``thermodynamic determinism.'' This objection assumes constraints should be negotiable---that sufficient cleverness or effort can circumvent physical limitations. The framework instead treats constraints as boundary conditions within which strategies operate. Constraints define the possibility space; strategies navigate within it. The Second Law does not prevent reconstruction; it specifies what reconstruction requires.

% ============================================
% END SECTION 9
% ============================================
%
% NEW CITATIONS REQUIRED FOR v2.1:
%
% @article{macnamara2014,
%   author = {Macnamara, Brooke N. and Hambrick, David Z. and Oswald, Frederick L.},
%   title = {Deliberate Practice and Performance in Music, Games, Sports, Education, 
%            and Professions: A Meta-Analysis},
%   journal = {Psychological Science},
%   volume = {25},
%   number = {8},
%   pages = {1608--1618},
%   year = {2014},
%   doi = {10.1177/0956797614535810}
% }
%
% @article{march1991,
%   author = {March, James G.},
%   title = {Exploration and Exploitation in Organizational Learning},
%   journal = {Organization Science},
%   volume = {2},
%   number = {1},
%   pages = {71--87},
%   year = {1991}
% }
%
% @article{frey2017,
%   author = {Frey, Carl Benedikt and Osborne, Michael A.},
%   title = {The Future of Employment: How Susceptible Are Jobs to Computerisation?},
%   journal = {Technological Forecasting and Social Change},
%   volume = {114},
%   pages = {254--280},
%   year = {2017}
% }
%
% @article{autor2015,
%   author = {Autor, David H.},
%   title = {Why Are There Still So Many Jobs? The History and Future of Workplace Automation},
%   journal = {Journal of Economic Perspectives},
%   volume = {29},
%   number = {3},
%   pages = {3--30},
%   year = {2015}
% }
%
% @book{sennett2008,
%   author = {Sennett, Richard},
%   title = {The Craftsman},
%   publisher = {Yale University Press},
%   year = {2008}
% }
%
% @book{stengers2018,
%   author = {Stengers, Isabelle},
%   title = {Another Science is Possible: A Manifesto for Slow Science},
%   publisher = {Polity Press},
%   year = {2018}
% }
%
% @misc{cotoaga2024karten,
%   author = {Cotoaga, Kurt},
%   title = {Karten auf den Tisch: Ein Framework f{\"u}r KI-informierte Entscheidungsfindung},
%   year = {2024},
%   url = {http://cotoaga.ai/karten-auf-den-tisch/},
%   note = {Accessed December 2024}
% }
%
% EXISTING CITATIONS VERIFIED:
% - korzybski1933, kahneman2011, hutchins1995, hofstadter1979
% - schein1985, dimaggio1983, granovetter1985, bourdieu1990
% - maturana1987, snowden2007, ford2024
% - DeLaCroix2018, collins2010, zuboff2019
% - mark2008, gonzalez2004 (for interruption studies)
%
% KEY CHANGES v2 -> v2.1:
% 1. Equation now references Section 4 methodology (no competing formulation)
% 2. Added "Operationalizing the Proxies" subsubsection with definitions
% 3. Added March (1991) for exploration/exploitation 20% slack anchor
% 4. Added Frey & Osborne, Autor for substitutability construct
% 5. Softened Chaotic domain: "non-authoritative" not "noise"
% 6. Added Section 9.5 "Boundary Conditions and Failure of Applicability"
% 7. Reframed existence proof with autoethnographic method note
% 8. Added external triangulation (Sennett, Stengers, complexity networks)
% 9. Fixed cross-reference: sec:equity not sec:measurement
% 10. De-escalated final paragraph (objection/response structure, not accusation)
% 11. Added table footnote linking evaluation signals to Section 9.8
% 12. Qualified "AI capability thresholds" -> "labor-market exposure intensifies"
% ============================================